\section{Introduction}
\label{sec:introduction}

Quantum algorithms for linear algebra and signal processing have emerged as a promising application domain for quantum computing, with quantum signal processing (QSP)~\cite{GSLW2019} and quantum singular value transformation (QSVT)~\cite{GSLW2019} providing a unified framework for manipulating quantum states in a controlled manner. Central to these techniques is the ability to encode classical information into quantum states and to perform transformations that implement desired polynomial functions. The choice of encoding strategy profoundly affects the expressiveness, efficiency, and geometric structure of quantum algorithms.

Traditional approaches to encoding classical data into quantum states include amplitude encoding, basis encoding, and angle encoding. While these methods have proven useful, they often impose constraints on the range of encodable values or lack natural geometric interpretations that could inform algorithm design. In particular, amplitude encoding restricts values to lie within the unit hypersphere, and basis encoding requires discretization of continuous variables.

\subsection{Motivation and Contributions}

In this work, we introduce \emph{stereographic encoding}, a novel quantum encoding strategy based on the stereographic projection from the complex plane $\C$ to the Bloch sphere $\Sphere$. The stereographic projection is a conformal mapping that has deep connections to complex analysis, projective geometry, and the Riemann sphere. By leveraging this classical geometric construction, we obtain a quantum encoding with several remarkable properties:

\begin{enumerate}
    \item \textbf{Unbounded encoding}: Unlike amplitude encoding, stereographic encoding can represent arbitrary complex numbers $z \in \C$ without magnitude constraints, as the encoding maps $\C \cup \{\infty\}$ bijectively to the Bloch sphere.
    
    \item \textbf{Geometric structure}: Single-qubit quantum gates induce Möbius (fractional linear) transformations on the encoded complex values, providing a rich algebraic and geometric structure. The Clifford group, in particular, generates automorphisms of the Riemann sphere.
    
    \item \textbf{Natural QSP connection}: Under stereographic encoding, the quantum signal processing protocol naturally generates Chebyshev polynomials and rational function approximations, with explicit formulas for zeros and poles.
    
    \item \textbf{Hamiltonian interpretation}: The encoding can be realized as the ground state of a parametrized Hamiltonian, enabling time-dependent control of trajectories in the complex plane.
\end{enumerate}

Our main contributions are:

\begin{itemize}
    \item We define the stereographic encoding $\ket{z} = \frac{1}{\sqrt{|z|^2+1}}(z\ket{0} + \ket{1})$ and derive its properties, including density matrix representation, Bloch vector components, and measurement-based decoding.
    
    \item We establish the connection between quantum gates and Möbius transformations, showing that gates in $\SU(2)$ act as fractional linear transformations on encoded complex values.
    
    \item We analyze quantum signal processing with stereographic encoding, deriving the Chebyshev polynomial structure of iterated signal operators and computing the resulting rational polynomial approximations for various degrees.
    
    \item We develop the Hamiltonian formulation, deriving the time-dependent driving terms needed to realize arbitrary trajectories in $\C$ and computing the geometric (Berry) phase accumulated along closed paths.
    
    \item We explore applications to quantum eigenvalue transformation, quantum kernel methods, and quantum classification.
\end{itemize}

\subsection{Organization}

The remainder of this paper is organized as follows. Section~\ref{sec:preliminaries} reviews the mathematical prerequisites, including stereographic projection and the Bloch sphere representation of qubits. Section~\ref{sec:encoding} introduces stereographic encoding and derives its fundamental properties. Section~\ref{sec:mobius} establishes the correspondence between quantum gates and Möbius transformations. Section~\ref{sec:qsp} analyzes quantum signal processing under this encoding, including the Chebyshev polynomial structure. Section~\ref{sec:dynamics} develops the Hamiltonian dynamics and Berry phase calculations. Section~\ref{sec:applications} discusses applications. Section~\ref{sec:discussion} concludes with a discussion of our results and future directions.
